\chapter{系统需求分析与总体设计}

\section{需求分析}

本系统面向布控球自适应 PTZ 跟踪场景，目标是在真实设备视频流上完成实时人体检测、目标选择、云台控制和可视化调试。与离线图像检测或单一算法验证不同，本系统需要同时连接真实布控球设备、处理实时码流、运行检测模型、维护目标状态并下发云台控制指令，因此需求分析必须覆盖算法效果、设备接口和运行稳定性。结合任务书、开题报告和中期检查材料，系统需求可分为功能需求、性能需求和工程约束三类。

\subsection{功能需求}

系统需要实现以下功能：

\begin{enumerate}
    \item 设备接入与视频预览。系统能够适配海康威视 SDK 接口，完成设备初始化、登录认证、实时预览启动和码流解码，并提供稳定的实时画面输入。
    \item 人体目标检测。系统能够从实时视频帧中检测人体目标，输出目标框、中心点、置信度和面积等信息，并形成可用于控制决策的检测结果。
    \item 自动 PTZ 跟踪。系统能够根据检测目标相对画面中心的偏差，自动生成水平和垂直方向的云台控制指令，使目标逐步回到画面中心附近。
    \item 多目标身份记忆。系统能够为画面中的人体目标分配 ID，支持自动跟踪最大人体目标，也支持用户指定某一 ID 进行持续跟踪。
    \item 人工接管与调试。系统提供 PTZ 人工控制、检测状态管理、自动跟踪状态管理、倒挂模式适配、目标选择和检测结果反馈等交互能力。
    \item 资源释放与异常处理。系统退出时需要停止云台运动，释放预览句柄、播放库端口和 SDK 登录资源，避免设备连接残留。
\end{enumerate}

从系统运行流程看，上述功能共同构成一个闭环：系统完成设备接入和实时预览后，检测模块持续输出人体目标信息，目标选择模块根据自动模式或指定 ID 模式确定跟踪对象，控制模块根据目标偏差生成 PTZ 指令。当设备倒挂安装时，倒挂适配机制对画面方向和控制方向进行统一修正。该流程要求系统既能自动运行，也保留必要的人工干预入口。

\subsection{性能需求}

布控球跟踪系统具有较强的实时性要求。检测模块需要在较短时间内完成单帧推理，使目标坐标能够及时输入控制模块；PTZ 控制模块需要限制指令频率，避免机械结构过度抖动；视频预览需要保持基本流畅，以便用户观察设备状态。根据中期检查材料，系统初步目标为视频流延迟控制在 200 ms 以内，检测帧率达到实时跟踪需求，并在正常光照和无遮挡场景下实现稳定人体跟踪。

除实时性外，系统还需要具备一定的稳定性和可恢复性。检测模型偶发推理失败时不应导致整个程序退出，视频帧暂时读取失败时控制模块不应继续使用过期目标，用户关闭窗口时系统必须停止云台动作并释放 SDK 资源。对于原型系统而言，这些要求不一定体现为复杂的容错框架，但需要通过异常捕获、状态判断和退出清理机制落实到程序流程中。

\subsection{工程约束}

系统开发受到设备接口、平台兼容性和算力条件的共同约束。当前原型需要适配海康威视 HCNetSDK 与 PlayCtrl 等设备接口，并通过 Python 完成算法模块、视频链路和控制逻辑的快速集成。若后续迁移至嵌入式开发板，还需要遵循串口、GPIO 等接口的电气规范和数据传输要求\cite{texas_gpio_2016}。目标检测依赖 ultralytics、OpenCV、Pillow 和 NumPy 等视觉计算库。由于真实布控球设备存在机械响应时间，PTZ 控制不能简单地按每帧检测结果连续发送指令，需要在控制策略中考虑死区、冷却和脉冲时长等参数。

此外，工程约束还体现在跨平台依赖和模块解耦上。设备接口、检测模型、目标记忆和控制策略在运行时具有不同的资源需求和异常模式，因此系统设计需要降低各模块之间的耦合程度。本文在总体设计中将设备接口、检测逻辑和目标记忆尽可能分离，使核心算法模块具备独立测试和后续迁移的可能。

\section{总体架构设计}

系统采用分层架构，整体可划分为视频采集层、检测分析层、控制决策层和用户交互层，如图\ref{fig:system-architecture}所示。

\begin{figure}[H]
    \centering
    \begin{tikzpicture}[
        node distance=0.8cm,
        layer/.style={draw, rounded corners=2pt, minimum width=11cm, minimum height=0.95cm, align=center},
        arrow/.style={->, thick}
    ]
        \node[layer] (ui) {用户交互层：状态管理、目标选择、人工接管、倒挂模式};
        \node[layer, below=of ui] (control) {控制决策层：目标选择、偏差计算、平滑滤波、PTZ 指令生成};
        \node[layer, below=of control] (detect) {检测分析层：YOLOv8n 人体检测、置信度过滤、人体 ReID 记忆};
        \node[layer, below=of detect] (video) {视频采集层：HCNetSDK 登录预览、PlayCtrl 解码、实时帧缓存};
        \draw[arrow] (video) -- (detect);
        \draw[arrow] (detect) -- (control);
        \draw[arrow] (control) -- (ui);
        \draw[arrow] (ui.east) -- ++(1.0,0) |- (control.east);
    \end{tikzpicture}
    \caption{系统总体架构}
    \label{fig:system-architecture}
\end{figure}

视频采集层负责将布控球实时码流转换为检测模块可读取的图像帧；检测分析层负责人体检测和目标身份记忆；控制决策层负责选择跟踪目标并生成云台控制指令；用户交互层用于展示视频、切换功能和进行手动干预。各层之间通过结构化数据传递，减少模块耦合。

该分层方式有助于降低系统复杂度。视频采集层只关心设备登录、预览和解码，不直接参与目标判断；检测分析层只输出候选目标和身份编号，不直接控制云台；控制决策层根据当前目标和系统状态生成动作；用户交互层只修改开关状态和显示结果。通过这种划分，单个模块发生异常时更容易定位问题，也便于后续替换检测模型或迁移到不同硬件平台。

\section{数据流设计}

系统运行时的数据流包括视频流、检测结果流和控制指令流。

\begin{enumerate}
    \item 视频流：布控球通过网络输出实时码流，HCNetSDK 的预览回调将码流数据送入 PlayCtrl，PlayCtrl 解码回调将最新帧保存为临时图像，并更新帧缓存信息。
    \item 检测结果流：YOLO 检测线程从帧缓存读取图像，执行人体检测，生成候选人体列表。候选列表经过 ReID 记忆模块更新后形成带有 person\_id 的目标列表。
    \item 控制指令流：目标选择模块根据自动模式或指定 ID 模式确定跟踪目标，PTZ 控制线程读取目标中心点，计算偏差并调用 NET\_DVR\_PTZControl\_Other 发送控制指令。
    \item 交互状态流：交互模块改变检测启停、自动跟踪、倒挂适配和指定目标 ID 等状态变量，从而影响检测、显示和控制逻辑。
\end{enumerate}

数据流设计的核心原则是使用最新状态而非积压队列。布控球跟踪关注当前目标位置，如果检测线程处理过旧帧，即使检测结果准确，也可能导致云台朝过期方向运动。因此，系统在帧缓存和检测结果上均采用“最新值覆盖”的方式，使控制模块尽可能读取最近一次有效检测结果。该策略牺牲了逐帧完整处理能力，但更符合实时控制场景对低延迟的要求。

同时，系统将交互状态作为独立状态流处理。检测状态、PTZ 自动跟踪状态、倒挂模式和目标选择模式发生变化时，系统并不直接修改底层设备连接，而是改变对应状态变量，由检测线程、显示线程和控制线程在各自循环中读取这些变量。这样可以避免交互事件直接执行耗时操作，从而保持系统响应。

\section{关键模块划分}

根据系统功能边界，关键模块及职责如表\ref{tab:module-design}所示。

\begin{table}[H]
    \centering
    \caption{系统关键模块划分}
    \label{tab:module-design}
    \begin{tabularx}{\textwidth}{c|X|X}
        \hline
        模块 & 输入输出 & 职责 \\
        \hline
        主控与交互模块 & 系统状态、用户指令 & 组织系统运行流程，管理功能状态、人工接管和资源释放 \\
        \hline
        视频采集模块 & 设备码流、图像帧 & 接收实时码流，完成解码和最新帧缓存，为预览和检测提供输入 \\
        \hline
        设备接口模块 & 控制指令、设备状态 & 封装设备登录、预览控制、码流回调和 PTZ 控制等底层接口 \\
        \hline
        人体 ReID 模块 & 检测框、外观特征 & 构造人体特征，维护短时目标记忆，分配 person\_id 并选择跟踪目标 \\
        \hline
        测试验证模块 & 模拟目标、运行日志 & 验证目标选择、ReID 合并、指定 ID 等核心逻辑 \\
        \hline
    \end{tabularx}
\end{table}

从模块划分可以看出，本文系统采用“主流程集中调度、关键逻辑独立封装”的实现方式。主控与交互模块负责连接界面、模型、设备和线程，是系统集成的核心；视频采集模块将实时预览和码流解码逻辑从主流程中分离，降低视频链路对其他模块的影响；人体 ReID 模块将目标记忆逻辑独立出来，使其能够通过单元测试验证。该组织方式虽然仍属于原型工程，但已经为后续进一步模块化提供了基础。

\section{线程模型设计}

系统采用多线程方式分离视频检测、PTZ 控制和检测结果反馈。主线程负责交互事件循环；YOLO 检测线程按固定间隔读取最新帧并更新检测结果；PTZ 控制线程周期性检查当前目标并执行云台控制；结果反馈线程绘制检测框、目标 ID、偏移量和跟踪状态。多线程设计能够避免耗时推理阻塞交互响应，也使控制线程能够以相对稳定的频率检查目标位置。

在线程协同方面，系统需要处理两个典型矛盾。第一，检测线程推理耗时相对较长，而界面显示和云台控制需要保持响应，因此检测不能放在主线程中执行。第二，PTZ 控制不能简单地在每次检测完成后立即触发，否则检测帧率变化会直接影响控制频率。本文将检测和控制拆分为两个循环，使控制线程能够基于最新检测结果独立判断是否需要动作，并通过冷却时间约束指令发送频率。

当前原型通过全局状态变量共享检测结果和控制开关。该方式实现简单，适合单机原型验证；后续若部署到长期运行设备，可引入线程安全队列、锁机制或消息总线，以进一步提高并发访问的可维护性。对于本文阶段而言，全局状态方式能够降低实现复杂度，使主要精力集中在闭环功能验证上；但在工程产品化阶段，应进一步明确状态读写边界，避免多线程竞争导致偶发错误。

\section{控制策略设计}

PTZ 控制的核心输入为目标中心坐标 $(x_t, y_t)$ 与画面中心坐标 $(x_c, y_c)$。系统定义偏差：

\begin{equation}
    e_x = x_t - x_c,\quad e_y = y_t - y_c
\end{equation}

当 $|e_x|$ 大于水平阈值时，系统根据 $e_x$ 的符号选择向左或向右转动；当 $|e_y|$ 大于垂直阈值时，系统根据 $e_y$ 的符号选择向上或向下转动。为降低检测框抖动对控制的影响，目标中心先经过指数平滑：

\begin{equation}
    \hat{x}_t^{(k)} = \alpha x_t^{(k)} + (1-\alpha)\hat{x}_t^{(k-1)}
\end{equation}

\begin{equation}
    \hat{y}_t^{(k)} = \alpha y_t^{(k)} + (1-\alpha)\hat{y}_t^{(k-1)}
\end{equation}

其中 $\alpha$ 为平滑系数。当前系统中 PTZ\_SMOOTH\_ALPHA 设置为 0.3，水平阈值 PTZ\_THRESHOLD\_X 为 60，垂直阈值 PTZ\_THRESHOLD\_Y 为 65。单次云台脉冲时长由基础时长和偏差比例项共同决定：

\begin{equation}
    T = \min(T_0 + \lambda(|e_x|+|e_y|), T_{\max})
\end{equation}

其中 $T_0$ 为基础脉冲时长，$\lambda$ 为比例系数，$T_{\max}$ 为上限。该策略能够在偏差较大时提高回正速度，在目标接近中心时减少过度转动。

除上述控制量外，系统还设置控制冷却时间，用于限制连续 PTZ 指令之间的最小间隔。冷却机制能够避免检测框在中心附近小幅波动时导致云台高频启停。对于倒挂安装场景，系统在计算控制方向前对画面和偏差方向进行统一处理，使用户从界面看到的运动方向与实际控制方向保持一致。总体而言，本文采用的是面向工程调试的轻量控制策略，参数含义清晰，便于根据不同设备和部署环境进行调整。
